\section{Risultati completi delle ablazioni controllate}\label{app:controlled-results}
Il disegno fattoriale comprende 64 ricette PSO e un riferimento DE sul campione F1, F8, F15, F21. Ogni funzione rappresenta rispettivamente separabilità, condizionamento moderato, multimodalità con struttura globale e multimodalità con struttura debole della classificazione BBOB. Il gruppo delle funzioni fortemente mal condizionate F10--F14 è assente. I ranghi per gruppo del CSV coincidono qui con una singola funzione per gruppo e dimensione.
\subsection{Struttura dei grafi e graduatoria}
\begin{table}[H]\centering\small
\caption{Proprietà delle realizzazioni a 64 nodi nel fattoriale, mediate sui 24 blocchi funzione--dimensione--istanza. Il grado esclude il self; il vicinato sociale lo include. Tutti i grafi risultano connessi.}\label{tab:graph-properties}
\begin{tabular}{@{}lrrr@{}}
\toprule Grafo & Grado medio & Cammino medio & Clustering\\\midrule
Global & 63.000 & 1.000 & 1.000\\
griglia & 4.000 & 4.063 & 0.000\\
R4 & 4.086 & 3.093 & 0.054\\
R8 & 8.202 & 2.178 & 0.134\\
ring & 2.000 & 16.254 & 0.000\\
scale-free & 3.906 & 2.778 & 0.196\\
small-world & 4.000 & 3.910 & 0.330\\
star & 1.969 & 1.969 & 0.000\\
\bottomrule\end{tabular}
\end{table}

\clearpage\begingroup\footnotesize\setlength{\tabcolsep}{6pt}\renewcommand{\arraystretch}{0.96}
\begin{longtable}{@{}llllrr@{}}
\caption{Graduatoria completa del fattoriale. A: propagatore analitico; D: aggiornamento discreto; S: casualità scalare; C: per coordinata. I ranghi sono attribuiti alle mediane di gap su tre istanze entro ciascuna delle quattro funzioni.}\label{tab:controlled-ranks}\\
\toprule Dinamica & Casualità & Grafo & $h$ & Rango $D=5$ & Rango $D=20$\\\midrule\endfirsthead
\toprule Dinamica & Casualità & Grafo & $h$ & Rango $D=5$ & Rango $D=20$\\\midrule\endhead
\bottomrule\endfoot
DE & --- & --- & --- & 15.50 & 9.50\\
A & C & Global & 0.38125 & 26.50 & 15.75\\
A & C & Global & 1 & 16.00 & 21.00\\
A & C & Global & 3.6 & 25.25 & 29.50\\
A & C & griglia & 0.38125 & 31.25 & 10.00\\
A & C & griglia & 1 & 11.75 & 3.25\\
A & C & griglia & 3.6 & 17.75 & 15.00\\
A & C & R4 & 0.38125 & 29.50 & 10.75\\
A & C & R4 & 1 & 10.00 & 4.75\\
A & C & R4 & 3.6 & 21.75 & 14.75\\
A & C & R8 & 0.38125 & 22.25 & 5.25\\
A & C & R8 & 1 & 10.75 & 1.75\\
A & C & R8 & 3.6 & 9.75 & 11.25\\
A & C & ring & 0.38125 & 41.75 & 41.00\\
A & C & ring & 1 & 14.25 & 17.75\\
A & C & ring & 3.6 & 40.50 & 36.25\\
A & C & scale-free & 0.38125 & 35.00 & 15.50\\
A & C & scale-free & 1 & 12.25 & 8.75\\
A & C & scale-free & 3.6 & 14.25 & 20.75\\
A & C & small-world & 0.38125 & 29.50 & 13.00\\
A & C & small-world & 1 & 11.75 & 6.25\\
A & C & small-world & 3.6 & 12.50 & 16.50\\
A & C & star & 0.38125 & 41.50 & 38.25\\
A & C & star & 1 & 30.50 & 29.50\\
A & C & star & 3.6 & 51.00 & 48.75\\
A & S & Global & 0.38125 & 34.50 & 27.50\\
A & S & Global & 1 & 28.25 & 40.50\\
A & S & Global & 3.6 & 39.75 & 55.25\\
A & S & griglia & 0.38125 & 26.25 & 35.75\\
A & S & griglia & 1 & 16.50 & 45.75\\
A & S & griglia & 3.6 & 33.50 & 60.00\\
A & S & R4 & 0.38125 & 26.50 & 39.00\\
A & S & R4 & 1 & 13.50 & 41.25\\
A & S & R4 & 3.6 & 42.50 & 55.75\\
A & S & R8 & 0.38125 & 20.75 & 29.25\\
A & S & R8 & 1 & 19.25 & 32.00\\
A & S & R8 & 3.6 & 32.50 & 54.25\\
A & S & ring & 0.38125 & 43.50 & 59.25\\
A & S & ring & 1 & 28.00 & 57.25\\
A & S & ring & 3.6 & 59.50 & 64.75\\
A & S & scale-free & 0.38125 & 29.75 & 41.00\\
A & S & scale-free & 1 & 19.25 & 46.00\\
A & S & scale-free & 3.6 & 54.75 & 60.25\\
A & S & small-world & 0.38125 & 20.50 & 37.50\\
A & S & small-world & 1 & 16.50 & 47.50\\
A & S & small-world & 3.6 & 48.75 & 62.00\\
A & S & star & 0.38125 & 36.25 & 54.50\\
A & S & star & 1 & 44.25 & 52.25\\
A & S & star & 3.6 & 52.75 & 64.25\\
D & C & Global & 1 & 47.50 & 31.25\\
D & C & griglia & 1 & 47.00 & 39.00\\
D & C & R4 & 1 & 50.50 & 37.75\\
D & C & R8 & 1 & 54.75 & 36.75\\
D & C & ring & 1 & 53.00 & 46.00\\
D & C & scale-free & 1 & 50.25 & 38.75\\
D & C & small-world & 1 & 55.00 & 34.25\\
D & C & star & 1 & 52.00 & 41.50\\
D & S & Global & 1 & 35.25 & 22.00\\
D & S & griglia & 1 & 39.50 & 25.25\\
D & S & R4 & 1 & 42.75 & 21.75\\
D & S & R8 & 1 & 37.25 & 23.25\\
D & S & ring & 1 & 51.50 & 45.00\\
D & S & scale-free & 1 & 46.00 & 32.00\\
D & S & small-world & 1 & 55.25 & 26.50\\
D & S & star & 1 & 57.25 & 36.50\\
\end{longtable}\endgroup

\subsection{Contrasti dei fattori}
L'effetto limitato è definito dall'Eq.~\eqref{eq:relative-effect}. A meno D confronta i propagatori allo stesso tipo di casualità e grafo, a $h=1$. C meno S confronta la casualità. R8 meno Global confronta i vicinati nella stessa dinamica e casualità. Le interazioni sono differenze di effetti limitati: per il passo si sottrae l'effetto R8--griglia a $h=1$ da quello all'altro passo; per la dinamica si sottrae l'effetto R8--Global discreto da quello analitico. Il segno delle interazioni indica il cambiamento del vantaggio relativo. Gli intervalli sono descrittivi del ricampionamento di quattro funzioni; il controllo di Holm comprende tutti gli 84 contrasti.
\small
\begin{longtable}{@{}rlllrrrr@{}}
\caption{Tutti gli 84 contrasti pianificati. L'effetto è la mediana fra funzioni degli effetti limitati appaiati; l'IC bootstrap ricampiona le quattro funzioni. Tutti i valori corretti sono uno.}\label{tab:planned}\\
\toprule $D$ & Contrasto & Tipo & Grafo/dinamica & Effetto & IC 95\% & $p$ & $p_{\mathrm{Holm}}$\\\midrule\endfirsthead
\toprule $D$ & Contrasto & Tipo & Grafo/dinamica & Effetto & IC 95\% & $p$ & $p_{\mathrm{Holm}}$\\\midrule\endhead
\bottomrule\endfoot
5 & A meno D & S & Global & -0.159 & [-1.000, 0.135] & 0.3750 & 1.0000\\
5 & A meno D & S & ring & -0.563 & [-0.993, -0.269] & 0.1250 & 1.0000\\
5 & A meno D & S & griglia & -0.606 & [-0.996, 0.318] & 0.2500 & 1.0000\\
5 & A meno D & S & star & -0.522 & [-0.882, 0.164] & 0.2500 & 1.0000\\
5 & A meno D & S & R4 & -0.638 & [-0.989, -0.047] & 0.1250 & 1.0000\\
5 & A meno D & S & R8 & -0.203 & [-0.999, -0.003] & 0.1250 & 1.0000\\
5 & A meno D & S & small-world & -0.617 & [-0.991, -0.239] & 0.1250 & 1.0000\\
5 & A meno D & S & scale-free & -0.359 & [-0.998, -0.196] & 0.1250 & 1.0000\\
5 & A meno D & C & Global & -0.560 & [-1.000, -0.005] & 0.1250 & 1.0000\\
5 & A meno D & C & ring & -0.646 & [-0.964, -0.565] & 0.1250 & 1.0000\\
5 & A meno D & C & griglia & -0.826 & [-1.000, -0.539] & 0.1250 & 1.0000\\
5 & A meno D & C & star & -0.593 & [-0.993, 0.360] & 0.2500 & 1.0000\\
5 & A meno D & C & R4 & -0.762 & [-0.999, -0.016] & 0.1250 & 1.0000\\
5 & A meno D & C & R8 & -0.619 & [-1.000, -0.616] & 0.1250 & 1.0000\\
5 & A meno D & C & small-world & -0.641 & [-1.000, -0.042] & 0.1250 & 1.0000\\
5 & A meno D & C & scale-free & -0.928 & [-1.000, -0.287] & 0.1250 & 1.0000\\
5 & R8 meno Global & S & analytic & 0.218 & [-0.185, 0.904] & 0.3750 & 1.0000\\
5 & R8 meno Global & C & analytic & 0.259 & [-0.351, 0.963] & 0.6250 & 1.0000\\
5 & R8 meno Global & S & discrete & 0.227 & [-0.312, 0.667] & 0.6250 & 1.0000\\
5 & R8 meno Global & C & discrete & 0.366 & [-0.084, 0.714] & 0.2500 & 1.0000\\
5 & C meno S & A & Global & -0.468 & [-0.712, -0.000] & 0.1250 & 1.0000\\
5 & C meno S & A & ring & -0.320 & [-0.942, -0.002] & 0.1250 & 1.0000\\
5 & C meno S & A & griglia & -0.488 & [-0.966, 0.228] & 0.3750 & 1.0000\\
5 & C meno S & A & star & -0.146 & [-0.939, 0.019] & 0.3750 & 1.0000\\
5 & C meno S & A & R4 & -0.045 & [-0.921, 0.586] & 0.6250 & 1.0000\\
5 & C meno S & A & R8 & -0.181 & [-0.894, 0.033] & 0.3750 & 1.0000\\
5 & C meno S & A & small-world & -0.092 & [-0.897, 0.252] & 0.6250 & 1.0000\\
5 & C meno S & A & scale-free & -0.519 & [-0.969, 0.135] & 0.3750 & 1.0000\\
5 & C meno S & D & Global & 0.263 & [0.029, 0.455] & 0.1250 & 1.0000\\
5 & C meno S & D & ring & -0.158 & [-0.397, 0.226] & 0.6250 & 1.0000\\
5 & C meno S & D & griglia & 0.230 & [-0.072, 0.580] & 0.2500 & 1.0000\\
5 & C meno S & D & star & -0.122 & [-0.474, 0.108] & 0.3750 & 1.0000\\
5 & C meno S & D & R4 & 0.138 & [-0.031, 0.523] & 0.2500 & 1.0000\\
5 & C meno S & D & R8 & 0.466 & [-0.315, 0.718] & 0.3750 & 1.0000\\
5 & C meno S & D & small-world & -0.012 & [-0.496, 0.246] & 0.8750 & 1.0000\\
5 & C meno S & D & scale-free & 0.305 & [0.074, 0.546] & 0.1250 & 1.0000\\
5 & Interazione $h=.38125$ & S & R8/griglia & -0.193 & [-0.362, 0.120] & 0.3750 & 1.0000\\
5 & Interazione $h=3.6$ & S & R8/griglia & 0.030 & [-0.104, 0.345] & 0.6250 & 1.0000\\
5 & Interazione dinamica & S & R8/Global & 0.184 & [-0.278, 0.309] & 0.6250 & 1.0000\\
5 & Interazione $h=.38125$ & C & R8/griglia & -0.040 & [-1.086, 0.592] & 0.8750 & 1.0000\\
5 & Interazione $h=3.6$ & C & R8/griglia & -0.110 & [-0.883, -0.049] & 0.1250 & 1.0000\\
5 & Interazione dinamica & C & R8/Global & 0.008 & [-0.684, 0.212] & 1.0000 & 1.0000\\
20 & A meno D & S & Global & 0.168 & [0.134, 0.950] & 0.1250 & 1.0000\\
20 & A meno D & S & ring & 0.199 & [0.155, 0.594] & 0.1250 & 1.0000\\
20 & A meno D & S & griglia & 0.507 & [0.098, 0.688] & 0.1250 & 1.0000\\
20 & A meno D & S & star & 0.403 & [0.057, 0.770] & 0.1250 & 1.0000\\
20 & A meno D & S & R4 & 0.331 & [0.141, 0.837] & 0.1250 & 1.0000\\
20 & A meno D & S & R8 & 0.094 & [0.054, 0.721] & 0.1250 & 1.0000\\
20 & A meno D & S & small-world & 0.390 & [0.045, 0.874] & 0.1250 & 1.0000\\
20 & A meno D & S & scale-free & 0.197 & [0.116, 0.834] & 0.1250 & 1.0000\\
20 & A meno D & C & Global & -0.256 & [-0.933, -0.146] & 0.1250 & 1.0000\\
20 & A meno D & C & ring & -0.631 & [-0.930, -0.153] & 0.1250 & 1.0000\\
20 & A meno D & C & griglia & -0.941 & [-1.000, -0.443] & 0.1250 & 1.0000\\
20 & A meno D & C & star & -0.603 & [-0.766, 0.262] & 0.2500 & 1.0000\\
20 & A meno D & C & R4 & -0.651 & [-1.000, -0.294] & 0.1250 & 1.0000\\
20 & A meno D & C & R8 & -0.918 & [-1.000, -0.428] & 0.1250 & 1.0000\\
20 & A meno D & C & small-world & -0.637 & [-0.999, -0.114] & 0.1250 & 1.0000\\
20 & A meno D & C & scale-free & -0.670 & [-0.993, 0.040] & 0.2500 & 1.0000\\
20 & R8 meno Global & S & analytic & -0.144 & [-0.229, 0.110] & 0.3750 & 1.0000\\
20 & R8 meno Global & C & analytic & -0.802 & [-0.996, -0.214] & 0.1250 & 1.0000\\
20 & R8 meno Global & S & discrete & 0.088 & [0.041, 0.697] & 0.1250 & 1.0000\\
20 & R8 meno Global & C & discrete & 0.072 & [-0.044, 0.748] & 0.3750 & 1.0000\\
20 & C meno S & A & Global & -0.267 & [-0.982, 0.101] & 0.2500 & 1.0000\\
20 & C meno S & A & ring & -0.730 & [-0.990, -0.380] & 0.1250 & 1.0000\\
20 & C meno S & A & griglia & -0.897 & [-1.000, -0.456] & 0.1250 & 1.0000\\
20 & C meno S & A & star & -0.549 & [-0.877, -0.156] & 0.1250 & 1.0000\\
20 & C meno S & A & R4 & -0.650 & [-1.000, -0.408] & 0.1250 & 1.0000\\
20 & C meno S & A & R8 & -0.722 & [-1.000, -0.417] & 0.1250 & 1.0000\\
20 & C meno S & A & small-world & -0.787 & [-1.000, -0.349] & 0.1250 & 1.0000\\
20 & C meno S & A & scale-free & -0.662 & [-0.999, -0.368] & 0.1250 & 1.0000\\
20 & C meno S & D & Global & 0.247 & [-0.054, 0.617] & 0.2500 & 1.0000\\
20 & C meno S & D & ring & -0.066 & [-0.319, 0.479] & 1.0000 & 1.0000\\
20 & C meno S & D & griglia & 0.200 & [0.146, 0.708] & 0.1250 & 1.0000\\
20 & C meno S & D & star & 0.184 & [-0.253, 0.483] & 0.6250 & 1.0000\\
20 & C meno S & D & R4 & 0.446 & [0.082, 0.671] & 0.1250 & 1.0000\\
20 & C meno S & D & R8 & 0.218 & [0.021, 0.733] & 0.1250 & 1.0000\\
20 & C meno S & D & small-world & 0.170 & [-0.095, 0.659] & 0.6250 & 1.0000\\
20 & C meno S & D & scale-free & 0.343 & [-0.378, 0.589] & 0.3750 & 1.0000\\
20 & Interazione $h=.38125$ & S & R8/griglia & 0.379 & [-0.203, 0.696] & 0.2500 & 1.0000\\
20 & Interazione $h=3.6$ & S & R8/griglia & 0.051 & [-0.340, 0.389] & 0.8750 & 1.0000\\
20 & Interazione dinamica & S & R8/Global & -0.154 & [-0.640, 0.003] & 0.2500 & 1.0000\\
20 & Interazione $h=.38125$ & C & R8/griglia & 0.175 & [-0.647, 0.483] & 1.0000 & 1.0000\\
20 & Interazione $h=3.6$ & C & R8/griglia & 0.114 & [-0.164, 0.510] & 0.6250 & 1.0000\\
20 & Interazione dinamica & C & R8/Global & -0.726 & [-1.444, -0.259] & 0.1250 & 1.0000\\
\end{longtable}\normalsize

La famiglia separata dei 128 contrasti tra CPSO R8 analitico scalare a $h=1$ e tutte le altre configurazioni è disponibile integralmente nei riepiloghi allegati, con vittorie, pareggi, sconfitte, effetto relativo, effetto biseriale e $p$ corretto. Anche in tale famiglia tutti i valori corretti sono uno.
\subsection{Target ed ERT}
Il riferimento $f^*$ è letto dall'Observer ufficiale. I target sono 21 soglie logaritmiche fra $10^2$ e $10^{-8}$; ciascuna dimensione contiene 12 problemi e 252 coppie problema--target. L'ERT per target somma il costo del primo raggiungimento per i successi e il costo finale della run per gli insuccessi, dividendo per i successi. La componente temporale usa i timestamp osservati nelle callback seriali e il tempo finale per gli insuccessi, secondo la medesima regola di censura. Il CSV conserva tutti i metodi e i target; la tabella mostra tre ricette rappresentative.
\begin{table}[H]\centering\small
\caption{Raggiungimenti di target selezionati nel campione controllato. Il trattino indica zero successi e ERT indefinito. Il CSV allegato riporta tutti i 21 target per tutti i metodi.}\label{tab:targets}
\begin{tabular}{@{}rllrrrr@{}}
\toprule $D$ & Metodo & Target & Successi & Run & ERT FE & ERT ms\\\midrule
5 & CPSO R8 S & 100 & 12 & 12 & 76.3 & 1.61\\
5 & CPSO R8 S & 1 & 4 & 12 & 5713.2 & 18.99\\
5 & CPSO R8 S & 0.01 & 3 & 12 & 8561.7 & 27.45\\
5 & CPSO R8 S & 0.0001 & 3 & 12 & 9608.3 & 29.62\\
5 & CPSO R8 C & 100 & 12 & 12 & 100.2 & 1.62\\
5 & CPSO R8 C & 1 & 6 & 12 & 3407.3 & 12.00\\
5 & CPSO R8 C & 0.01 & 4 & 12 & 6165.0 & 20.53\\
5 & CPSO R8 C & 0.0001 & 3 & 12 & 9443.7 & 30.15\\
5 & DE & 100 & 12 & 12 & 137.8 & 2.28\\
5 & DE & 1 & 7 & 12 & 2954.9 & 17.94\\
5 & DE & 0.01 & 4 & 12 & 6655.2 & 38.44\\
5 & DE & 0.0001 & 2 & 12 & 14995.0 & 84.64\\
20 & CPSO R8 S & 100 & 7 & 12 & 7866.6 & 35.83\\
20 & CPSO R8 S & 1 & 0 & 12 & --- & ---\\
20 & CPSO R8 S & 0.01 & 0 & 12 & --- & ---\\
20 & CPSO R8 S & 0.0001 & 0 & 12 & --- & ---\\
20 & CPSO R8 C & 100 & 12 & 12 & 2298.8 & 11.80\\
20 & CPSO R8 C & 1 & 4 & 12 & 22847.2 & 109.32\\
20 & CPSO R8 C & 0.01 & 3 & 12 & 35208.0 & 167.24\\
20 & CPSO R8 C & 0.0001 & 1 & 12 & 119942.0 & 552.76\\
20 & DE & 100 & 9 & 12 & 5232.7 & 38.51\\
20 & DE & 1 & 4 & 12 & 24544.0 & 157.62\\
20 & DE & 0.01 & 4 & 12 & 27823.0 & 173.44\\
20 & DE & 0.0001 & 1 & 12 & 122720.0 & 751.83\\
\bottomrule\end{tabular}
\end{table}

